\section{Problema}
Imagine que toma posesión de la lámpara mágica. Tras frotarla, el genio de la lámpara mágica le cuenta algo sobre las clases de complejidad y se lo demuestra matemáticamente. En cada numeral le cuenta algo distinto.

Utilice esa información, junto con todo lo que la humanidad sabe sobre ellas, para dibujar todos los diagramas de contenencia posibles entre las clases de complejidad \(\mathrm{P}\), \(\mathrm{NP}\), \(\mathrm{NP\text{-}completo}\), \(\mathrm{NP\text{-}hard}\) y \(\text{co-NP}\).

Represente cada clase mediante una región. Una región dentro de otra indica que la primera clase está contenida por completo en la segunda. Si dos clases son iguales, represéntelas usando una sola región etiquetada con ambos nombres. Puede haber uno, varios o ningún diagrama posible: dibújelos todos, cada uno marcado con la condición que lo hace válido. Si no hay ninguno, explique por qué.

\section{Entregables}

\begin{enumerate}
  \item El genio le dice que \(\mathrm{P} = \mathrm{NP}\), \(\mathrm{P} = \text{co-P}\) y \(\text{co-NP} = \mathrm{NP}\).
  \item El genio le dice que \(\mathrm{P} = \text{co-P}\), \(\mathrm{P} \neq \mathrm{NP}\) y \(\mathrm{NP} \cap \text{co-NP} = \varnothing\).
  \item El genio le dice que \(\text{co-NP} \cup \mathrm{NP} = \text{co-NP}\) y \(\mathrm{P} \neq \mathrm{NP}\).
\end{enumerate}
